\documentclass[a4paper]{scrartcl}
% \usepackage{tgheros}
% \usepackage[T1]{fontenc}
\usepackage{fontspec}
\setsansfont{TeX Gyre Heros}[Ligatures=TeX]
\usepackage[all]{nowidow}

\usepackage[sorting=none]{biblatex}
\addbibresource{references.bib}

\usepackage{graphicx}
\graphicspath{{./figure}}
\usepackage{trimclip}
% \usepackage{xspace}
\usepackage{csquotes}
\usepackage{hyperref}
\usepackage{todonotes} % [disable]
\usepackage{amsmath}
\usepackage{siunitx}
% \newfontfamily{\lmr}{Latin Modern Roman}
% \sisetup{text-font-command=\lmr}
% \sisetup{output-prefix={\ensuremath{\mu}}}
\usepackage{amssymb}

\newcommand{\figref}[1]{Fig.\,\ref{#1}}


\begin{document}

% \begin{titlepage}
\begin{center}
\vspace*{1cm}

\LARGE
  \textbf{Control and Regulation Systems}

\vspace{.2cm}

\normalsize
  Physical Measuring Techniques FK7063

\vspace{.5cm}
{Michael Vogt}

May 2026
\end{center}

\vspace{1cm}



\section*{Introduction}
In this laboratory exercise, a PID control system will be used to control a pump with the
goal of keeping a constant water level in a small tank. The control system has
three tweakable parameters which will be combined in different ways
to optimize the pump's behaviour.

% \end{titlepage}



\section{Theory and Background}
Control systems are used to regulate the behaviour of a system (the \textit{controlled system})
that depends on various variables. In the case of this experiment, the controlled system is
a pump leading into a water tank whose water level is measured by a sensor. The control system
is a microcontroller (Arduino) that receives information about the water level and can
regulate the speed of the pump.

PID controllers are a type of control system that regulates control variables (the speed of the pump)
based on the difference between the controlled variable (the water level in the tank) %$e(t) = r(t) - y(t)$
and a reference value (the desired water level). This difference is called the error $e(t)$
and a PID system regulates the control variable according to three terms derived from the error:
\begin{itemize}
  \item P -- Proportional $\proptop e(t)$
  \item I -- Integral $\propto \int e(t) \text{d}t$
  \item D -- Differential $\proptop \frac{\text{d}e}{\text{d}t}$
\end{itemize}
The strength of the individual terms can be controlled by three parameters $P$, $I$ and $D$.



\section{Setup and Experiment}\label{Setup}
The experimental setup consists of a pump leading into a water tank. The pump is controlled
by an Arduino based on the water level in the tank. This works by means of a PID controller
implemented in software. This implementation is accompanied by a graphical user interface
that can be used to change parameters and watch parameter changes over time.

Different parameter combinations are tested to understand the behaviour of the PID system.

\subsection{Proportional Term}
\begin{figure}
  \centering
  \includegraphics[width=.6\textwidth]{p1}
  \caption{Water level changes with only the propotional ($P$) term.
    Parameters and coefficient values are shown as sliders near the top of the window.
    The order is: target level, maximum level, $P$, $I$, $D$. The left two graphs
    show the water level over time.}
  \label{fig:p1}
\end{figure}

This is because the pressure acting back on the pump is higher with a higher water level,
making the same pump activation cause a lesser increase in the level.



\section{Modulation Experiment}



\clearpage
\section{Conclusion}



\vfill
\printbibliography

\end{document}